%%%%%%%%%%%%%% DOCUMENT CLASS & PACKAGES %%%%%%%%%%%%%%
\documentclass{ctexart}
\usepackage{amsmath}
\usepackage{algorithm}
\usepackage{geometry}
\usepackage{minted}
\usemintedstyle[verilog]{xcode}
\usepackage[dvipsnames]{xcolor}
\usepackage{fontspec}
\usepackage{graphicx}
\usepackage{tabularx}
\usepackage{lmodern}
\usepackage{setspace}
\usepackage[hidelinks, bookmarksnumbered=false]{hyperref}
\usepackage{bookmark}
\usepackage[normalem]{ulem}
\usepackage{paralist}
\usepackage{subcaption}

%%%%%%%%%%%%%% STYLE SETTINGS %%%%%%%%%%%%%%
% 字体设置
\setCJKfamilyfont{hwzs}{STZhongsong}
\setmonofont{Menlo}

% 代码样式设置
\setminted{fontfamily=tt,frame=single,fontsize=\small,framesep=4mm,baselinestretch=1.2,linenos,breaklines}

% 页面布局
\geometry{a4paper,left=1in,right=1in,top=1.3in,bottom=1.3in}
\setstretch{1.5}

%%%%%%%%%%%%%% TITLE PAGE INFO %%%%%%%%%%%%%%

% 生成封面页命令
\newcommand{\maketitlepage}{
\begin{spacing}{1.2}
\begin{titlepage}
    \begin{figure}[H]
        \centering
        \includegraphics[width=0.3\textwidth]{./zjulogo.png}
    \end{figure}
    \vspace{20ex}
    \begin{center}
    \Huge {\CJKfamily{hwzs}\reporttitle} \\ [1ex]
    \Large {\CJKfamily{hwzs}\reportsubtitle} \\ [24ex]

    \begin{tabularx}{0.6\textwidth}{@{}lX@{}}
        \Large \textit{姓名：} & \Large \textit{\underline{\makebox[\linewidth][c]{\authorname}}} \\ [1ex]
        \Large \textit{学号：} & \Large \textit{\underline{\makebox[\linewidth][c]{\studentid}}} \\ [1ex]
        \Large \textit{邮箱：} & \Large \textit{\underline{\makebox[\linewidth][c]{\authoremail}}} \\ [1ex]
        \Large \textit{教师：} & \Large \textit{\underline{\makebox[\linewidth][c]{\teachername}}} \\ [1ex]
        \Large \textit{助教：} & \Large \textit{\underline{\makebox[\linewidth][c]{\assistantnames}}} \\ [1ex]
        \Large \textit{日期：} & \Large \textit{\underline{\makebox[\linewidth][c]{\today}}}
    \end{tabularx}
    \end{center}
\end{titlepage}
\end{spacing}
}

%%%%%%%%%%%%%% SECTION FORMATTING %%%%%%%%%%%%%%
% 设置计数器格式
\renewcommand{\thesection}{\chinese{section}}
\renewcommand{\thesubsection}{\chinese{subsection}、}

% 定义实验标题格式
\newcommand{\LabName}[1]{
    {\centering
    \section*{#1}
    \par}
    \addcontentsline{toc}{section}{#1}
    \markboth{#1}{#1}
}

% 定义小节标题格式
\newcommand{\LabSection}[1]{
    \stepcounter{subsection}
    \subsection*{\thesubsection #1}
    \addcontentsline{toc}{subsection}{\thesubsection #1}
}

% 设置section样式
\CTEXsetup[format={\Large\bfseries\centering}]{section}